\documentclass[12pt]{article}
\usepackage[utf8]{inputenc}
\usepackage[spanish]{babel}
\usepackage{amsmath}
\usepackage{amssymb}
\usepackage{booktabs}
\usepackage{geometry}
\geometry{margin=2.5cm}

\title{\textbf{Resumen: Hendry (1980)}\\
\large ``Econometrics---Alchemy or Science?''\\
\normalsize \textit{Economica}, Vol.~47, No.~188, pp.~387--406}
\author{David F. Hendry\\The London School of Economics}
\date{}

\begin{document}
\maketitle
\tableofcontents
\newpage

%----------------------------------------------------------------------
\section{Alquimia y Ciencia (\textit{Alchemy and Science})}
%----------------------------------------------------------------------

\subsection{Connotaciones de ``Ciencia''}

La \textbf{ciencia / science} es un proceso p\'ublico. Usa sistemas de conceptos llamados \textbf{teor\'ias / theories} para interpretar y unificar observaciones denominadas \textbf{datos / data}; a su vez, los datos se emplean para contrastar las teor\'ias. La creaci\'on de teor\'ias puede ser inductiva, pero su demostraci\'on y contrastaci\'on son deductivas; en materias inexactas, la contrastaci\'on involucra \textbf{inferencia estad\'istica / statistical inference}.

Las teor\'ias que son a la vez simples, generales y coherentes son las m\'as valiosas porque facilitan una pr\'actica cient\'ifica productiva y precisa. Una mayor restrictividad aumenta los riesgos de rechazo emp\'irico y con ello la \textbf{plausibilidad / plausibility} si la teor\'ia no es refutada. En la pr\'actica, los cambios de opini\'on cient\'ifica son lentos; seg\'un Kuhn (1962), el progreso cient\'ifico opera mediante cambios ``\textbf{revolucionarios / revolutionary}'' en los marcos te\'oricos b\'asicos.

\subsection{Connotaciones de ``Alquimia''}

La \textbf{alquimia / alchemy} denota el arte presunto de transmutar metales en metales nobles. Las connotaciones p\'eyorativas de la alquimia se asocian al oscurantismo, al misticismo y a ``recetas'' para simular oro que pod\'ian enga\~nar a los propios alquimistas. La relevancia de estos comentarios para el estado actual de la econometr\'ia se hace evidente m\'as adelante en el texto.

%----------------------------------------------------------------------
\section{Econometr\'ia (\textit{Econometrics})}
%----------------------------------------------------------------------

\subsection{Definici\'on y Alcance}

La \textbf{Sociedad Econom\'etrica / Econometric Society}, fundada en 1930, defini\'o su objetivo principal como ``el avance de la teor\'ia econ\'omica en su relaci\'on con la estad\'istica y las matem\'aticas''. En sentido amplio, la econometr\'ia analiza las relaciones entre variables econ\'omicas abstrayendo los fen\'omenos principales de inter\'es y enunciando teor\'ias en forma matem\'atica; su utilidad emp\'irica se eval\'ua usando informaci\'on estad\'istica.

En sentido estricto (utilizado en el art\'iculo), la econometr\'ia concierne la interpretaci\'on y an\'alisis de los datos en el contexto de la teor\'ia econ\'omica ``establecida''. Como expres\'o Frisch (1933): ``la penetraci\'on mutua de la teor\'ia econ\'omica cuantitativa y la observaci\'on estad\'istica es la esencia de la econometr\'ia''.

La teor\'ia econom\'etrica es el estudio de:
\begin{itemize}
  \item Las propiedades de los \textbf{procesos generadores de datos / data generation processes}.
  \item Las t\'ecnicas para analizar datos econ\'omicos.
  \item Los m\'etodos para estimar magnitudes num\'ericas de par\'ametros.
  \item Los procedimientos para contrastar hip\'otesis econ\'omicas.
\end{itemize}

\subsection{Cr\'iticas Hist\'oricas a la Econometr\'ia}

Desde Keynes (1939, 1940), numerosos autores han cuestionado el car\'acter cient\'ifico de la econometr\'ia. Entre las cr\'iticas m\'as citadas figuran las de Worswick (1972), quien sugiri\'o que algunos econometr\'istas producen meras \textbf{``herramientas de apariencia'' / ``pretend-tools''}, y Phelps Brown (1972), quien afirm\'o que las regresiones entre series temporales son ``susceptibles de enga\~nar''. Leontief (1971) calific\'o la econometr\'ia de ``intento de compensar la debilidad notoria de la base de datos mediante el uso m\'as y m\'as sofisticado de t\'ecnicas estad\'isticas''.

El autor responde a estas cr\'iticas demostrando la \textbf{alquimia emp\'irica / empirical alchemy} en acci\'on, con el fin de entender en qu\'e sentido las cr\'iticas son v\'alidas y c\'omo pueden sugerirse estrategias constructivas para reforzar el papel del m\'etodo cient\'ifico en la econometr\'ia.

%----------------------------------------------------------------------
\section{La Econometr\'ia como Alquimia (\textit{Econometrics as Alchemy})}
%----------------------------------------------------------------------

\subsection{La ``Piedra Filosofal'' Econom\'etrica}

La ``Piedra Filosofal'' de los econometristas es el \textbf{an\'alisis de regresi\'on / regression analysis}: una herramienta capaz de transformar datos en resultados ``significativos''. El enga\~no se practica f\'acilmente a partir de ``recetas'' falsas destinadas a simular hallazgos \'utiles; tales ejercicios son denominados despec\-tivamente \textbf{``regresiones sin sentido'' / ``nonsense regressions''} o \textbf{``regresiones espurias'' / ``spurious regressions''} en la profesi\'on.

\subsection{Fuentes de Resultados Alqu\'imicos}

El argumento central de esta secci\'on es que \textit{antes} de realizar estas regresiones, la teor\'ia relevante permiti\'o al autor \textit{deducir} qu\'e ocurrir\'ia y c\'omo construir los ejemplos deseados desde el primer intento. Esta capacidad predictiva previa es la que distingue el conocimiento cient\'ifico genuino de la mera alquimia.

Dos contrastes fundamentales revelan si un modelo es alqu\'imico o cient\'ifico:
\begin{enumerate}
  \item \textbf{Contraste de predicci\'on / prediction test}: un modelo alqu\'imico predice mal fuera de la muestra.
  \item \textbf{Contraste de constancia de par\'ametros / parameter constancy test}: un modelo alqu\'imico exhibe par\'ametros inestables a lo largo del tiempo.
\end{enumerate}

\subsection{Requisito Esencial para Modelos No Experimentales}

En una disciplina \textbf{no experimental / non-experimental}, un requisito esencial para que un modelo sea \'util es que pueda explicar por qu\'e modelos \textit{falsos} previos produjeron sus resultados observados. Esta propiedad es denominada \textbf{abarcamiento / encompassing} (Davidson \textit{et al.}, 1978).

%----------------------------------------------------------------------
\section{Problemas de la Econometr\'ia (\textit{Econometrics' Problems})}
%----------------------------------------------------------------------

\subsection{Lista de Problemas del Modelo de Regresi\'on Lineal}

Keynes (1939) identific\'o una lista de problemas del \textbf{modelo de regresi\'on lineal / linear regression model}, que en la terminolog\'ia moderna incluyen:

\begin{itemize}
  \item \textbf{Sesgo por variables omitidas / omitted variables bias}: conjunto incompleto de factores determinantes.
  \item \textbf{Variables no observables / unobservable variables}: por ejemplo, las expectativas.
  \item \textbf{Datos mal medidos / badly measured data}: estimados a partir de n\'umeros \'indice inapropiados.
  \item \textbf{Correlaciones espurias / spurious correlations}: derivadas del uso de variables \textit{proxy} y de la simultaneidad.
  \item \textbf{Multicolinealidad / multicollinearity}: incapacidad de separar los efectos de variables altamente correlacionadas.
  \item \textbf{Formas funcionales incorrectas / incorrect functional forms}: suponer linealidad cuando la relaci\'on es no lineal.
  \item \textbf{Din\'amica mal especificada / mis-specified dynamics}: longitudes de retardo e interacciones din\'amicas incorrectas.
  \item \textbf{Pre-filtrado incorrecto de datos / incorrect pre-filtering}: transformaciones previas que distorsionan la relaci\'on.
  \item \textbf{Inferencia causal inv\'alida / invalid causal inference}: deducir causalidades a partir de correlaciones.
  \item \textbf{Par\'ametros no constantes / non-constant parameters}: predicci\'on imprecisa por inestabilidad.
  \item \textbf{Significancia econ\'omica vs. estad\'istica / economic vs. statistical significance}: confundir ambos conceptos.
\end{itemize}

A esta lista el autor a\~nade: mala especificaci\'on estoc\'astica, supuestos de \textbf{exogeneidad / exogeneity} incorrectos, muestras insuficientes, \textbf{agregaci\'on / aggregation}, falta de \textbf{identificaci\'on estructural / structural identification} e incapacidad de retrotraer un\'ivocamente los resultados emp\'iricos a una teor\'ia inicial.

\subsection{Brecha entre Teor\'ia y Pr\'actica Emp\'irica}

La pr\'actica emp\'irica ha tendido a rezagarse con respecto a la ``frontera'' te\'orica. Los grandes sistemas macroeconometr\'icos estaban seriamente mal especificados en su estructura de \textbf{correlaci\'on serial / correlation structure}, lo que determin\'o su fracaso predictivo durante la crisis del petr\'oleo.

\subsection{Problemas de los Datos Econ\'omicos}

Los datos econ\'omicos son notooriamente poco fiables (Morgenstern, 1950). Variables como la ``renta personal real disponible'' son extremadamente dif\'iciles de medir con precisi\'on. Las discrepancias en la medici\'on del ingreso pueden tener importantes implicaciones de pol\'itica econ\'omica.

%----------------------------------------------------------------------
\section{Una Estructura para la Econometr\'ia (\textit{A Structure for Econometrics})}
%----------------------------------------------------------------------

\subsection{El Sistema Econ\'omico como Proceso Complejo}

El sistema econ\'omico es el resultado de siglos de comportamiento humano adaptativo. Los econometristas lo conceptualizan como un proceso din\'amico \textbf{no lineal, interdependiente, multivariante y de desequilibrio / nonlinear, interdependent, multivariate, disequilibrium dynamical process}, sujeto a perturbaciones aleatorias e involucrando muchos fen\'omenos no observables. Esta conceptualizaci\'on es la base real de la cr\'itica de Keynes.

\subsection{Estructura Formal: Proceso Generador de Datos}

Como primera aproximaci\'on, el \textbf{proceso generador de datos / data generation process (DGP)} en econom\'ia puede concebirse como:

\begin{equation}
  \mathbf{y}_t \mid \mathbf{z}_t \;\sim\; N(\mathbf{\Pi z}_t,\, \mathbf{\Omega}), \qquad t = 1, \ldots, T,
\end{equation}

donde $\mathbf{y}_t$ es un vector de variables end\'ogenas, $\mathbf{z}_t$ es un vector de toda la informaci\'on pasada y presente relevante, $E(\mathbf{y}_t/\mathbf{z}_t) = \mathbf{\Pi z}_t$, y $\mathbf{x}_t \sim NI(\mu, \Psi)$. La matriz de par\'ametros $(\mathbf{\Pi}, \mathbf{\Omega}) = \mathbf{P}$ se toma como aproximadamente constante.

\subsection{Teor\'ia Econ\'omica como Restricci\'on}

Una ``teor\'ia econ\'omica'' corresponde a afirmar que $\mathbf{P}$ depende solo de un n\'umero menor de par\'ametros, recogidos en el vector $\boldsymbol{\theta}$:

\begin{equation}
  \mathbf{P} = \mathbf{f}(\boldsymbol{\theta}), \qquad \boldsymbol{\theta} \in \Theta.
\end{equation}

Si $\boldsymbol{\theta}$ es \textbf{identificable / identifiable} (es decir, implicado un\'ivocamente por $\mathbf{P}$), entonces todas las hip\'otesis del tipo~(2) pueden contrastarse usando el principio de Wald (1943). La funci\'on de verosimilitud es:

\begin{equation}
  L(\mathbf{P};\, \mathbf{y}_t \mid \mathbf{z}_t).
\end{equation}

\subsection{Exogeneidad D\'ebil y Factorizaci\'on de la Verosimilitud}

Si $L(.)$ puede factorizarse de forma que:

\begin{equation}
  L(.) = L_1(\theta_1,\, \mathbf{y}_{1t}/\mathbf{y}_{2t},\, \mathbf{z}_t)\cdot
         L_2(\theta_2;\, \mathbf{y}_{2t}/\mathbf{z}_t),
\end{equation}

de modo que cambios en $\theta_1$ no afectan a $\theta_2$ (par\'ametros ``nuisance''), entonces $L_1(.)$ puede analizarse de forma separada de $L_2(.)$ (Koopmans, 1950). En tal caso, se dice que $\mathbf{y}_{2t}$ es \textbf{d\'ebilmente ex\'ogena / weakly exogenous} para $\theta_1$ y puede tomarse como dada al analizar el submodelo que determina $\mathbf{y}_{1t}$.

\subsection{Ecuaci\'on Generadora de Estimadores}

La \textbf{estimaci\'on / estimation} concierne las reglas alternativas para asignar valores num\'ericos a $\theta_1$ dados los datos. Tomando $L = \log_e L$, la condici\'on de primer orden de la verosimilitud es:

\begin{equation}
  \frac{\partial L_1}{\partial \theta_1} = \mathbf{q}_1(\theta_1),
\end{equation}

denominada \textbf{ecuaci\'on generadora de estimadores / estimator-generating equation}, pues otros estimadores pueden interpretarse como \textit{aproximaciones} a la soluci\'on de $\mathbf{q}_1(\theta_1) = \mathbf{0}$ (Hendry, 1976). La inferencia es en gran medida dependiente de $\mathbf{q}(.)$.

\subsection{Modelizaci\'on Econom\'etrica}

La \textbf{modelizaci\'on econom\'etrica / econometric modelling} es la actividad de ajustar una versi\'on incorrecta de~(2) a una representaci\'on inadecuada de~(1). El compromiso resultante puede ser torpe o puede ser una aproximaci\'on \'util que recoja resultados previos, es suficientemente constante para la predicci\'on y la pol\'itica, y arroja luz sobre la teor\'ia econ\'omica. Escribir simplemente una ``teor\'ia econ\'omica'' y calibrar sus par\'ametros con un estimador pseudo sofisticado bas\'andose en datos pobres que el modelo no describe adecuadamente es una receta para el desastre: no para simular oro sino para el autoenga\~no.

\subsection{Supuesto ``ad hoc'' y Programas de Investigaci\'on}

El enfoque del autor es admitidamente \textbf{\textit{ad hoc}}, dado que la ``optimizaci\'on'' es un principio organizador sensato para la teor\'ia econ\'omica solo si las funciones de criterio asociadas representan adecuadamente los problemas de decisi\'on de los agentes (objetivos, costes, restricciones). El programa de investigaci\'on emp\'irico del autor se basa en supuestos \textbf{m\'inimos / minimal} sobre la inteligencia de los agentes, con m\'axima dependencia de los datos usando directrices de la ``teor\'ia econ\'omica'' para restringir la clase de modelos considerados.

%----------------------------------------------------------------------
\section{¿Es la Econometr\'ia Alquimia o Ciencia? (\textit{Is Econometrics Alchemy or Science?})}
%----------------------------------------------------------------------

\subsection{Testabilidad como Criterio Cient\'ifico}

La facilidad con que se pueden crear resultados espurios sugiere alquimia, pero el car\'acter cient\'ifico de la econometr\'ia queda ilustrado por el hecho de que tales enga\~nos son \textit{contrastables}. En un mundo que cambia r\'apidamente, las falacias no detectadas se convierten r\'apidamente en ejemplos positivos de la \textbf{``Ley'' de Goodhart / Goodhart's ``Law''} (1978), seg\'un la cual todos los modelos econom\'etricos se deterioran cuando se usan para la pol\'itica econ\'omica.

\subsection{El Estatus Cient\'ifico de la Econometr\'ia}

Es dif\'icil proporcionar un argumento convincente contra la acusaci\'on de Keynes de que la econometr\'ia es \textbf{``alquimia estad\'istica'' / ``statistical alchemy''}, puesto que muchas de sus cr\'iticas siguen siendo pertinentes. Sin embargo, la caracterizaci\'on de la ciencia ofrecida en la Secci\'on~I no excluye \textit{a priori} a la econometr\'ia simplemente por su incapacidad de realizar experimentos controlados. La sustanciaci\'on emp\'irica del car\'acter cient\'ifico requiere la existencia de \textbf{evidencia creible / credible evidence}: hallazgos aceptables con independencia de las creencias pol\'iticas o las preconcepciones sobre la estructura de la econom\'ia.

La turbulencia de la d\'ecada de 1970 facilit\'o enormemente el rechazo de modelos ``falsos'', y aunque todav\'ia estamos lejos de producir ``respuestas'' definitivas, se han logrado avances notables desde Keynes, aunque al coste de hacer el tema altamente t\'ecnico y cada vez m\'as inaccesible para los no especialistas.

\subsection{Las Tres Reglas de Oro de la Econometr\'ia}

Las \textbf{tres reglas de oro de la econometr\'ia / three golden rules of econometrics} son: \textbf{contrastar, contrastar y contrastar / test, test and test}. Aunque las tres reglas son violadas regularmente en las aplicaciones emp\'iricas, esto es f\'acilmente remediable. Los modelos rigurosamente contrastados que:
\begin{enumerate}
  \item describen adecuadamente los datos disponibles,
  \item \textbf{abarcan / encompass} los hallazgos previos, y
  \item se derivan de teor\'ias bien fundamentadas
\end{enumerate}
mejorar\'ian en gran medida cualquier pretensi\'on de car\'acter cient\'ifico de la econometr\'ia.

\subsection{Conclusi\'on}

Si la econometr\'ia resultar\'a ser m\'as an\'aloga a la alquimia que a la ciencia depende principalmente del esp\'iritu con que se aborde la materia. Hendry concluye que, aunque los avances son reales, el progreso m\'as r\'apido se lograr\'ia si todos los estudios emp\'iricos proporcionaran informaci\'on de contraste mejorada que permita a los lectores juzgar correctamente la plausibilidad de los modelos.

\end{document}
